\documentclass[
  12pt,			% Schriftgröße
	a4paper,			% Papierformat
	oneside,			% einseitiges Dokument
	titlepage,			% es wird eine Titelseite verwendet
	parskip=full,		% Abstand zwischen Absätzen (volle Zeile)=full halbe zeille = half
	headings=normal,	% Größe der Überschriften verkleinern
	BCOR=10mm,		% Bindekorrektur
	final				% Status des Dokuments (final/draft)
]{scrreprt}

\usepackage{graphicx}
%\usepackage{wasysym}
\usepackage{pstricks,pst-node,pst-tree} %Für Bäume
\usepackage{framed}  %Fuer die Hinweisbox
\usepackage{soul} %Fuer farbige Unterstreichungen
\usepackage{color} %Fuer schreibbare Farben wie z. B. red

%%Smiley ;-)
\newcommand{\smiley}{
\begin{pspicture}(0.5,0.35)
	\pscircle(0.25,0.15){0.2}
	\qdisk(0.19,0.20){0.03}
	\qdisk(0.31,0.20){0.03}
	\psarc(0.25,0.15){0.1}{210}{330}
\end{pspicture}
}

\newcommand{\blitz}{
\begin{pspicture}(0.25,0.25)
	\psscalebox{0.5}{
		\psline[linewidth=1pt](0.35,0.5)(0.15,0.25)
		\psline[linewidth=0.5pt](0.15,0.25)(0.35,0.25)
		\psline[linewidth=1pt](0.35,0.25)(0,0)
		\psline[linewidth=1pt](0.15,0)(0,0)
		\psline[linewidth=1pt](0,0.15)(0,0)
	}
\end{pspicture}
}

% Hinweisbox
\newenvironment{balken}{
  \def\FrameCommand##1{\textcolor{red}{\vrule width 3pt} \hspace{10pt}##1}
  \MakeFramed {\advance\hsize-\width \FrameRestore}}
 {\endMakeFramed}
\newcommand*{\hinweis}{\hspace{-15pt}\colorbox{red}{\textbf{\textcolor{white}{ HINWEIS }}}} 

\usepackage[latin1]{inputenc}
\usepackage[T1]{fontenc}
\usepackage{setspace}
\usepackage[
	automark,		% Kapitelangaben in Kopfzeile automatisch erstellen
	headsepline,		% Trennlinie unter Kopfzeile
	ilines			% Trennlinie linksbündig ausrichten
]{scrpage2}
\usepackage{lmodern}	%Schriftbild in Adobe PDf sonst unleserlich
 
%Formatierung geht ab hier los:
 
\setkomafont{sectioning}{\normalcolor\bfseries} % Ändert Überschriftschriftart (nun mit Serifen)
\onehalfspacing																	%Zeilenabstand 1,5 fach
\KOMAoptions{DIV=last}		%Satzspiegel neu berechnen wegen 1,5 fach Zeilenabstand
 
% Kopf- und Fußzeilen ------------------------------------------------------
\pagestyle{scrheadings}
 
% Kopf- und Fußzeile auch auf Kapitelanfangsseiten -------------------------
\renewcommand*{\chapterpagestyle}{scrheadings}
 
% Schriftform der Kopfzeile ------------------------------------------------
\renewcommand{\headfont}{\normalfont}
 
% Kopfzeile ----------------------------------------------------------------
\ihead{\textit{\headmark}}
\chead{}
\chead{}
\setlength{\headheight}{21mm} % Höhe der Kopfzeile
\setheadwidth[0pt]{textwithmarginpar} % Kopfzeile über den Text hinaus verbreitern
\setheadsepline[text]{0.4pt} % Trennlinie unter Kopfzeile
 
% Fußzeile -----------------------------------------------------------------
%\ifoot{\copyright\ \autor}
\cfoot{}
\ofoot{\pagemark}

\title{Automatentheorie und formale Sprachen}
\subject{Script zur Lehrveranstaltung}
\author{Florian Schleich}
\subtitle{Prof. Annika Wagner}
\publishers{\includegraphics[scale=0.5]{images/HS-Fulda-logo.eps}}

\begin{document}
\maketitle

\tableofcontents

\chapter{Grammatiken}

In diesem Kapitel lernen wir, wie man Grammatiken definiert und W�rter aus diesen Grammatiken erzeugt.

\section{Symbole}
Symbole sind alle verwendbaren Zeichen um durch Kombination W�rter zu erschaffen.

\section{Alphabet}
Ein Alphabet, abgek�rzt $\Sigma$ (Sigma), ist eine endliche Menge von einzelnen Symbolen. Ein Alphabet wird als endliche Menge von Symbolen definiert.

$\Sigma = \{\smiley, a, Anton\}$

Neben einzelnen Buchstaben k�nnen auch W�rter im Sinne der deutschen Sprache verwendet werden. Es k�nnen aber auch andere Symbole (\smiley) verwendet werden. Bei Verwendung von W�rtern als Symbole eines Alphabets ist aber zu beachten, dass bei der sp�teren Verwendung die einzelnen Buchstaben oder Teile des Wortes nicht als Symbole genommen werden d�rfen, sondern das gesamte Wort als ein Symbol bezeichnet wird. Im hier verwendeten Fall darf also das gro�e $A$ nicht als Symbol verwendet werden darf. ebenso d�rfen auch keine Wortteile wie $ton$ nicht verwendet werden. Ein zul�ssiges Symbol ist allerdings das gesamte Wort $Anton$.

\section{W�rter}
W�rter werden aus Kombination von einzelnen Symbolen gebildet. Die bildbaren W�rter werden in einer Menge $\Sigma^{*}$ (Sigma mit Kleene-Stern) gesammelt. Wird die Menge aller W�rter gebildet muss jetzt explizit das leere Wort $\epsilon$ (Epsilon) angegeben werden. Auf Grund der in der Grammatik verwendeten Regeln wird die Menge der m�glichen W�rter eingeschr�nkt. In unserem Beispiel kann die Menge aller W�rter so aussehen:

$\Sigma^{*} = \{\epsilon, \smiley, a, Anton, \smiley a, a Anton, \smiley Anton, a \smiley, Anton a, Anton \smiley, \smiley \smiley, ...\}$

Die Menge wird gebildet, in dem erst das leere Wort, dann alle einstelligen W�rter, alle zweistelligen W�rter, alle dreistelligen W�rter usw. in die Menge integriert werden. Die Menge der W�rter ist unendlich, da die L�nge der W�rter nicht bestimmt ist.

Die L�nge eines Wortes $|w|$ bezeichnet die Anzahl der f�r das Wort verwendeten Symbole.

\begin{figure}[h]
	$|\epsilon| = 0$
	
	$|a| = 1$
	
	$|Anton| = 1$
	
	$|aa| = 2$
	\caption{Beispiele f�r Wortl�ngen}
\end{figure}

\section{Grammatiken}
Eine Grammatik $G$ wird definiert durch:
\begin{itemize}
	\item eine Menge von Variablen $V=\{S,U\}$, die bei der Ableitung verschwinden,
	\item eine Menge von Terminalsymbolen $\Sigma$, gelichbedeutend mit dem definierten Alphabet
	\item ein Startsymbol $S \in V$
	\item und eine Menge von Regeln $P$
\end{itemize}

Formal wird eine Grammatik also beschrieben als:

$G = (V, \Sigma, S, P)$

\subsection{Regeln}
Eine Regel besteht immer aus einer linken und einer rechten Regelseite (Tupel). Bei der Ableitung von Regeln wird der Ausdruck auf der \textcolor[rgb]{0,0,1}{linken Seite} durch den Ausdruck auf der \textcolor[rgb]{0,0.58,0}{rechten Seite} ersetzt.

Bei der Regel $\textcolor[rgb]{0,0,1}{S} \rightarrow \textcolor[rgb]{0,0.58,0}{Anton}$ wird die Variable $S$ durch das Symbol $Anton$ ersetzt. 

\subsection{Ableitung}
Bei der Ableitung von Regeln beginnt man beim Startsymbol (hier $S$) und ersetzt das Startsymbol durch den entsprechenden Ausdruck einer Regel, die auf das Startsymbol passt. Zwischen zwei Ableitungsschritten kommt dann ein Doppelpfeil ($\Rightarrow$). Eine k�rzere Schreibweise erlaubt der Doppelpfeil mit Kleene-Stern ($\Rightarrow^{*}$), bei dem nicht jede Regel explizit aufgeschrieben werden muss. Auch 0-Schritte sind bei dieser Ableitungsart zul�ssig (siehe auch Abschnitt \ref{ref-trans-sigma-kleene}, Seite \pageref{ref-trans-sigma-kleene}).

\begin{figure}[h]
$S \rightarrow Anton$ (Regel der Grammatik)

$S \Rightarrow Anton$ (Ableitung, direkte Anwendung der Regel ohne Zwischenschritte)

$S \Rightarrow^{*} Anton$ (Ableitung, mehrmaliges Anwenden von Regeln mit Zwischenschritten)

\caption{Einfaches Beispiel f�r eine Ableitung}
\end{figure}

\subsection{Regelalternativen}
Regeln k�nnen auch Alternativen enthalten, um mehr M�glichkeiten zu bieten W�rter zu erzeugen.

Alternativen werden auf der rechten Regelseite durch ein Pipe-Symbol ($|$) eingeleitet.
Eine Regel mit Alternativen k�nnte folgenderma�en aussehen:

$S \rightarrow Anton | Anton \smiley | Anton a$

Dadurch w�ren die W�rter $Anton$, $Anton \smiley$ und $Anton a$ f�r die Variable $S$ m�glich.

\subsection{Rekursion in Regeln}
\begin{figure}[h]
	$S \rightarrow Anton G$
	
	$G \rightarrow \smiley | \smiley G$
	\caption{Grammatik mit Rekursion}
\end{figure}

Rekursionen in Grammatiken k�nnen W�rter erzeugen, deren L�nge (Betrag) unendlich lang sind. Eine m�gliche Ableitung f�r die obige Grammatik k�nnte folgenderma�en lauten.

\begin{figure}[h]
	$S \Rightarrow Anton G$
	
	$\Rightarrow Anton \smiley G$
	
	$\Rightarrow Anton \smiley \smiley$
	\caption{Ableitung f�r eine rekursive Grammatik}
\end{figure}

Wir sehen also, dass durch die Rekursion auf $G$ die W�rter unendlich lang werden k�nnen und durch das Symbol \smiley aufgef�llt werden. Des weitern stellen wir fest, dass Ableitungen erst enden, wenn alle Variablen entfernt, bzw. ersetzt wurden.

\subsection{Mengen von Regeln}
Als weitere Komponente f�r die Definition von Grammatiken wird die Menge aller Regeln der Grammatik ben�tigt. Im vorherigen Abschnitt haben wir gelernt, wie Regeln definiert werden. Die Menge der Regeln einer Grammatik ist eine Teilmenge des kartesische Produktes der Elemente aus $V$ und $\Sigma$. Das kartesische Produkt wird durch die Kombination aller Elemente \textcolor[rgb]{0,0,1}{einer Menge} (alles vor dem $\times$) mit s�mtlichen Elementen einer \textcolor[rgb]{0,0.58,0}{anderen Menge} (alles nach dem $\times$) gebildet.

\begin{figure}[h]
	$A = \{0,1\}$
	
	$B = \{x\}$
	
	$\textcolor[rgb]{0,0,1}{A} \times \textcolor[rgb]{0,0.58,0}{B} = \{(\textcolor[rgb]{0,0,1}{0},\textcolor[rgb]{0,0.58,0}{x}),(\textcolor[rgb]{0,0,1}{1},\textcolor[rgb]{0,0.58,0}{x})\}$
	\caption{Kartesisches Produkt der Mengen A und B}
\end{figure}

Die Menge aller Regeln einer Grammatik wird nun durch $(V\cup\Sigma)^*\times(V\cup\Sigma)^*$ gebildet. Durch diese Definition lassen sich s�mtliche Kombinationen aus Terminalsymbolen und Variablen auf beden Seiten der Regeln ausfstellen. Es ist aber meistens gew�nscht nur ein paar bestimmte Regeln zuzulassen. Man bildet also eine bestimmte Teilmenge der eben gebildeten Gesamtmenge.

\begin{balken}
\hinweis{}
\textbf{Unterschied zwischen $\{\}$ und $()$} \newline
Bei Verwendung von $\{\}$ ist die Reihenfolge und Anordnung der Elemente egal, da es sich um eine Menge handelt.
Bei Verwendung von $()$ ist die Anordnung der Elemente \textbf{nicht} egal.
\end{balken} 

\subsection{�berpr�fung der Regeln}
Um sicherzustellen, dass eine Grammatik eine bestimmte Sprache erzeugt, m�ssen �berpr�fungen bez�glich der Vollst�ndigkeit und Korrektheit dieser Grammatik erfolgen (siehe auch Abschnitt \ref{alt-notation-grammatik}, Seite \pageref{alt-notation-grammatik}).

\paragraph{Vollst�ndigkeit}
Eine Grammatik ist vollst�ndig, wenn sie (alle) W�rter der Sprache erzeugt.

\paragraph{Korrektheit}
Eine Grammatik ist korrekt, wenn sie keine W�rter ausserhalb der Sprache erzeugt.

\section{Kontextfrei oder regul�r}
Kontextfreie Grammatiken haben auf der linken Regelseite genau \textbf{eine} Variable und eine Beliebige Kombination aus Terminalsymbolen und Variablen auf der rechten Regelseite. Grammatiken sind regul�r, wenn die verwendeten Regeln auf der linken Regelseite \textbf{eine} Variable und auf der rechten Seite entweder das leere Wort ($\epsilon$), genau ein Terminalsymbol oder genau ein Terminalsymbol und eine Variable besitzen.

\begin{figure}[h]
\begin{tabbing}
    \hspace*{1cm}\=\hspace*{3cm}\=\hspace*{3cm}\=\hspace*{2.7cm}\= \kill
    \= \textbf{Allgemein:} \> $P \subseteq (V \cup \Sigma)^{*} \times (V \cup \Sigma)^{*}$ \\
    \= \textbf{Kontextfrei:} \> $P \subseteq (V) \times (V \cup \Sigma)^{*}$ \\
		\= \textbf{Regul�r:} \> $P \subseteq (V) \times (Y \cup \Sigma \cup \{\epsilon\})$ wobei $Y = aV$
\end{tabbing}
\end{figure}

\begin{balken}
\hinweis{}
\textbf{Links- und Rechtsregul�r} \newline
In dieser LVA wurde der Begriff "regul�r" von Frau Wagner als "rechtsregul�r" (immer zuerst rechteste Variable ersetzen) definiert!
\end{balken} 

\section{Alternative Schreibweise f�r Grammatiken}\label{alt-notation-grammatik}
Um Grammatiken und deren Regeln einfacher kommunizieren zu k�nnen kann man die erzeugten W�rter auch als "Formel" aufschreiben. Dazu folgende Aufgabe.

\textbf{Aufgabe:} Sei $\Sigma = \{0,1\}$. Leiten Sie mit der folgenden Grammatik das Wort 0101001010 ab.

\begin{figure}[h]
$S \rightarrow 0T0$

$T \rightarrow 1S1 | \epsilon$
\end{figure}

\textit{Ableitung:}
\begin{figure}[h]
	$S \Rightarrow 0T0$
	
	$\Rightarrow \underline{0}1S1\underline{0}$
	
	$\Rightarrow \underline{01}0T0\underline{10}$
	
	$\Rightarrow \underline{010}1S1\underline{010}$
	
	$\Rightarrow \underline{0101}0T0\underline{1010}$
	
	$\Rightarrow \underline{01010}\underline{01010}$
	\caption{Zugeh�rige Ableitung}
\end{figure}

Aus dieser Grammatik lassen sich folgende m�gliche W�rter ableiten:

\begin{itemize}
	\item 00
	\item 010010
	\item 0101001010
	\item ...
\end{itemize}

Da die Grammatik rekursiv ist und damit die Wortliste nicht endlich ist, kann man unm�glich alle W�rter auflisten. Um alle W�rter die erzeugt werden k�nnen trotzdem aufschrieben zu k�nnen, kann die Menge der W�rter auch formel�hnlich aufgeschrieben werden.

$\{(01)^n00(10)^n|n\geq0\}$

Der erste Teil (01) ist genau so oft enthalten wie der letzte Teil (10) kann aber auch weggelassen werden. der mittlere Teil (00) ist immer vorhanden und kann auch nicht wiederholt werden. Durch diese vereinfachte Schreibweise ist es sehr leicht m�glich weitere W�rter zu erzeugen. Man muss nur einen Wert f�r $n$ einsetzen und wiederholt die entsprechenden Teile $n$-mal.

\textbf{Aufgabe:} Sei $\Sigma=\{0,1\}$. Entwicklen Sie eine Grammatik, die folgende Sprache erzeugt: $L=\{w \in \Sigma^{*} | \left|w\right| \;ist\;gerade\}$

\textit{Ansatz:} Das das zu erzeugende Wort ($w$) ein $\Sigma^{*}$ enthalten sein soll, kann es also aus allen m�glichen Kombinationen aus $\Sigma$ (0 und 1) bestehen. Als zus�tzliche Beschr�nkung f�r die Menge aller W�rter soll der Betrag, also die L�nge, des Wortes gerade (durch Zwei teilbar) sein.

Aus dem obigen Ansatz ergeben sich folgende Regeln f�r die Grammatik.

\begin{figure}[h]
\begin{tabbing}
    \hspace*{1cm}\=\hspace*{4.2cm}\=\hspace*{3cm}\=\hspace*{2.7cm}\= \kill
    \= $S \rightarrow PS | P$ \> (beliebig viele P�rchen aneinanderreihen) \\
    \= $P \rightarrow ZZ | \epsilon$ \> (P�rchen von Zahlen aus Ziffern bilden) \\
		\= $Z \rightarrow 0|1$ \> (Ziffern)
\end{tabbing}
\end{figure}

\begin{figure}[ht]
$S \Rightarrow PS$

$\Rightarrow \underline{P}P$

$\Rightarrow ZZ\underline{P}$

$\Rightarrow \underline{ZZ}ZZ$

$\Rightarrow 0\underline{ZZZ}$

$\Rightarrow \underline{0}1\underline{ZZ}$

$\Rightarrow \underline{01}0\underline{Z}$

$\Rightarrow \underline{0100}$
\caption{M�gliche Ableitung f�r das Wort 0100}
\end{figure}

\section{Ableitungsbaum}
Aus dieser Ableitung l�sst sich ein Ableitungsbaum erstellen. Dieser stellt strukturiert dar, wie ein Wort abgeleitet wird.

\begin{figure}[h]
\pstree[levelsep=8ex]{\Tcircle{S} }{
	\pstree{ \Tcircle{P} }{
		\pstree{ \Tcircle{Z} }{
			\Tcircle{0}
		}
		\pstree{ \Tcircle{Z} }{
			\Tcircle{1}
		}
	}
	\pstree{ \Tcircle{S} }{
		\pstree{ \Tcircle{P} }{
			\pstree{ \Tcircle{Z} }{
				\Tcircle{0}
			}
			\pstree{ \Tcircle{Z} }{
				\Tcircle{0}
			}
		}
	}
}
\caption{Ableitungsbaum f�r die obige Ableitung}
\end{figure}

Egal in welcher Reihenfolge abgeleitet wird, bleibt der Ableitungsbaum immer der selbe. Die Ableitung nach $S \Rightarrow PS \Rightarrow ZZS \Rightarrow ...$ oder nach $S \Rightarrow PS \Rightarrow PP \Rightarrow ...$ erzeugen den selben Ableitungsbaum. Anders sieht es aus, wenn eine andere Regelalternative zur Ableitung genutzt wird. Dann kann sich der Ableitungsbaum durchaus unterscheiden.

Die Ableitung $S \Rightarrow PS \Rightarrow PPS \Rightarrow PPP \Rightarrow ZZPP \Rightarrow ZZZZP \Rightarrow ZZZZ\epsilon \Rightarrow ...$ erzeugt den folgenden Ableitungsbaum.

\begin{figure}[ht]
\pstree[levelsep=8ex]{\Tcircle{S} }{
	\pstree{ \Tcircle{P} }{
		\pstree{ \Tcircle{Z} }{
			\Tcircle{0}
		}
		\pstree{ \Tcircle{Z} }{
			\Tcircle{1}
		}
	}
	\pstree{ \Tcircle{S} }{
		\pstree{ \Tcircle{P} }{
			\pstree{ \Tcircle{Z} }{
				\Tcircle{0}
			}
			\pstree{ \Tcircle{Z} }{
				\Tcircle{0}
			}
		}
		\pstree{ \Tcircle{S} }{
			\pstree{ \Tcircle{P} }{
				\Tcircle{$\epsilon$}
			}
		}
	}
}
\caption{Ableitungsbaum f�r eine alternative Ableitung}
\end{figure}

\begin{balken}
\hinweis{}
\textbf{Mehrdeutigkeit} \newline
Da es f�r die hier verwendete Grammatik strukturell unterschiedliche Ableitungsb�ume gibt, ist die Grammatik mehrdeutig. Daher ist es nicht m�glich einen Parser f�r die erzeugte Sprache zu generieren.
\end{balken}

Durch Linksableitung kann Mehrdeutigkeit von B�umen (Ableitungen) eindeutig bewiesen werden. Die Linksableitung ist eine Konvention, dass von Ableitungsschritt zu Ableitungsschritt die linkeste Variable zuerst abgeleitet wird.

\subsection{Reflexiver \& transitiver Abschluss von $\Rightarrow$} \label{ref-trans-sigma-kleene}
Der Ableitungsoperator $\Rightarrow$ erlaubt es, nur einen Ableitungsschritt darzustellen. Der Ableitungsoperator $\Rightarrow^*$ hingegen erlaubt es beliebg viele Ableitungsschritte zusammenzufassen. Wird eine Grammatik mit dem einfachen Operator so abgeleitet $S \Rightarrow bB \Rightarrow bb$, kann die Ableitung mit dem reflexiven und trasitiven Abschluss so aussehen: $S \Rightarrow^* bb$. Dies spart Zeit und schreibarbeit, ist aber nicht so �bersichtlich, da zwischeninformationen �ber die Ableitungen verloren gehen.

\section{Greibach Normalform (GNF)}
Die Greibach Normalform erzeugt bei jedem ableitungsschritt genau ein Terminalsymbol. Bei der GNF besteht die rechte Regelseite aus genau einem Terminalsymbol und beliebig vielen Variablen ($P \subseteq V \times \Sigma V^{*}$).

Im folgenden werden Regeln besprochen, die aus einer kontextfreien Grammatik eine Grammatik in Greibach Normalform erzeugen.

\subsection{Problem 1: Mehrere Terminalsymbole}
Wenn eine Regel mehrere Terminalsymbole auf der rechten Regelseite besitzt, verst��t sie gegen die GNF. Um das Problem zu l�sen, erzeugt man einfach eine neue Variable, die das Terminalsymbol ersetzt.

\begin{figure}[h]
\begin{tabbing}
    \hspace*{0cm}\=\hspace*{3cm}\=\hspace*{3cm}\=\hspace*{3cm}\= \kill
    \= $S \rightarrow a\underline{a}$ \> wird zu \> $S \rightarrow a\underline{A}$ \\
		\= \> \> $\underline{S \rightarrow a}$
\end{tabbing}
\end{figure}

\subsection{Problem 2: Mehrere Variablen, aber kein Terminalsymbol}
Ein weiteres Problem entsteht, wenn zwar mehrere Variablen auf der rechten Regelseite stehen, aber kein Terminalsymbol vorhanden ist. Dazu wird die vorderste Variable durch die rechte Regelseite der Variable ersetzt.

\begin{figure}[h]
\begin{tabbing}
    \hspace*{0cm}\=\hspace*{3cm}\=\hspace*{3cm}\=\hspace*{3cm}\= \kill
    \= $S \rightarrow \underline{B}S | c$ \> wird zu \> $S \rightarrow \underline{b}S | c$ \\
		\= $B \rightarrow b$ \> \> $B \rightarrow b$
\end{tabbing}
\end{figure}

\subsection{Problem 3: Epsilon}
Bei der GNF ist die Verwendung von $\epsilon$ generell untersagt. �berall wo $\epsilon$ einsetzbar ist, muss eine neue rechte Seite ohne $\epsilon$ entstehen (Variable die zu $\epsilon$ f�hrt muss also gestrichen werden).

\begin{figure}[h]
\begin{tabbing}
    \hspace*{0cm}\=\hspace*{3cm}\=\hspace*{3cm}\=\hspace*{3cm}\= \kill
    \= $S \rightarrow aA$ \> wird zu \> $S \rightarrow \underline{a} | aA$ \\
		\= $A \rightarrow bA | \underline{\epsilon}$ \> \> $A \rightarrow bA | \underline{b}$
\end{tabbing}
\end{figure}

Dabei ist zu beachten, dass das k�rzeste Wort \textbf{pro Regel} erhalten bleiben muss (f�r $A$ wurde jeweils $\epsilon$ eingesetzt).

$S \Rightarrow aA \Rightarrow \underline{a}$ \newline
$A \Rightarrow bA \Rightarrow \underline{b}$
\chapter{Endliche Automaten}
Endliche Automaten werden dazu genutzt um zu pr�fen ob ein Wort auf eine bestimmte Grammatik passt. Diese Eigenschaft nutzt ein Parser um den Quellcode auf seine syntaktische Richtigkeit zu pr�fen.
\end{document}